% Front/abstract.tex
% -----------------------------------------------------------------------------
% The abstract is an unnumbered chapter (*) so it does not get a chapter number,
% but we still add it to the table of contents manually with \addcontentsline.
% NTU specifies a maximum of one page for the abstract.
%
% A good abstract for an engineering FYP covers four things in order:
%   1. The problem and why it matters (1--2 sentences)
%   2. What you did / your approach (2--3 sentences)
%   3. Your key results with numbers (2--3 sentences)
%   4. The significance / conclusion (1 sentence)

\chapter*{Abstract}
\addcontentsline{toc}{chapter}{Abstract}

Large language models (LLMs) such as TinyLlama have demonstrated strong
natural language understanding capabilities but require substantial
computational resources that limit their deployment in embedded and
edge environments. Field-programmable gate arrays (FPGAs) offer a
compelling alternative to GPU-based inference by providing reconfigurable
hardware with low power consumption and predictable latency.

This project implements and optimises the inference pipeline of TinyLlama-1.1B
on the Xilinx ZCU102 FPGA platform using Vitis High-Level Synthesis (HLS).
The core transformer components -- including RMSNorm, matrix--vector
multiplication, multi-head attention, and the SwiGLU feed-forward network --
were implemented as quantised INT8 HLS kernels with hardware-level pipeline
and loop unrolling optimisations to minimise the initiation interval.

% --- Replace this block with your actual results once available ---------------
% Example of the kind of quantitative results you should report here:
%   "The final implementation achieves a throughput of X tokens/second at
%    Y MHz clock frequency, utilising 43.9\% of available LUTs and 34.2\%
%    of BRAM on the ZCU102. This represents a Z$\times$ improvement in
%    energy efficiency compared to a software baseline running on an
%    ARM Cortex-A53 processor."
% -----------------------------------------------------------------------------

The results demonstrate that FPGA acceleration with INT8 quantisation
achieves competitive inference performance while operating within the
ZCU102's resource budget, validating the viability of transformer-based
language model deployment on reconfigurable hardware.
